\documentclass[11pt]{article}
\usepackage{amsmath,amssymb,amsthm}
\usepackage[margin=1in]{geometry}
\usepackage{booktabs,longtable,array}
\usepackage[colorlinks=true,linkcolor=blue,urlcolor=blue]{hyperref}

\newtheorem{lemma}{Lemma}
\newtheorem{proposition}{Proposition}
\theoremstyle{definition}
\newtheorem{definition}{Definition}
\theoremstyle{remark}
\newtheorem{remark}{Remark}

\newcommand{\Tref}{\mathcal{T}}
\newcommand{\eps}{\epsilon}
\newcommand{\nhat}{\hat n}
\newcommand{\dd}{\mathrm{d}}
\newcommand{\sky}[1]{\langle #1\rangle_{\Omega}}
\newcommand{\ens}[1]{\langle #1\rangle}
\newcommand{\status}[1]{\textsf{[#1]}}

\title{Proposed approach: analytical temperature power spectrum\\
through fourth order in a radiation-only universe}
\author{Draft for review}
\date{2026-10-31}

\begin{document}
\maketitle

\begin{abstract}
We propose an exact reformulation of the problem that removes most of the
bookkeeping that usually makes higher-order CMB calculations error-prone.
The reformulation uses the fluid's own velocity potential and temperature as
coordinates and conformal factor. For an irrotational radiation fluid, the
observed temperature then becomes an \emph{exact} single line integral of
the shear of a conformally rescaled synchronous metric:
\[
T_{\rm obs}(\nhat)=T_{\rm fluid}(o)\,
\exp\Big[-\!\int_{t_e}^{t_o}\!\sigma_{ij}e^ie^j\,\dd t\Big].
\]
The observer's local temperature cancels exactly in the normalised map, so
endpoint, observer, Doppler, Sachs--Wolfe and integrated Sachs--Wolfe effects
are no longer separate pieces that must be combined consistently. They come
out automatically from one expression, evaluated at events fixed by physical
labels. The plan has five stages: exact lemmas; the perturbative solution
through $O(\eps^3)$ with nonlinear Gaussian $\zeta$ initial data; expansion
of the observable; angular and ensemble reduction to $C_\ell^{(11)}$,
$C_\ell^{(22)}$ and $C_\ell^{(13)}$; and monochromatic specialisation. It
comes with a validation matrix in which every check is exact symbolic
algebra.
Stage~0 (exact lemmas) and Stage~1 (complete first order, including the
closed-form transfer function) are \emph{implemented and verified} in
\texttt{verification/}. The remaining stages are specified here in enough
detail to execute. The open items are listed explicitly in
\S\ref{sec:open}.
\end{abstract}

\tableofcontents

%======================================================================
\section{Reading of the problem and scope}
%======================================================================

The problem asks for an end-to-end calculation that has, to our knowledge,
not been completed in the literature for any cosmology. The full
third-order map $\Theta^{(3)}$ including all photon, endpoint, observer and
normalisation terms, and hence the one-loop $C_\ell$, is new. Second-order
CMB results exist (Bartolo--Matarrese--Riotto; Pitrou--Uzan--Bernardeau;
Huang--Vernizzi; Fidler et al.; the SONG code). They are, however, built on
Thomson-scattering Boltzmann hierarchies and numerics, which are irrelevant
or forbidden here. The problem as posed is simpler in two decisive ways that
we exploit:
\begin{enumerate}
\item the test photons are collisionless from emission onward, so ``Boltzmann''
reduces to the null-geodesic equation;
\item the matter is a single barotropic, conformally invariant perfect fluid
($P=\rho/3$, $\rho=\alpha T^4$), which admits an exact velocity-potential
description.
\end{enumerate}
It is harder in another: all three orders must be analytic, with finite
times and every boundary term kept.

\paragraph{Deliverable structure.} The final answer is
\[
C_\ell=\eps^2C_\ell^{(11)}+\eps^4\Big[\tfrac14 C_\ell^{(22)}+\tfrac13\operatorname{Re}C_\ell^{(13)}\Big],
\]
for $\ell\ge2$. The definitions of $C_\ell$ for $\ell=0,1$ are discussed in
\S\ref{sec:dipole}. Each $C^{(rs)}_\ell$ is to be given as an explicit
functional of the primordial spectrum. It is a finite-dimensional definite
integral with closed-form (elementary, $\operatorname{Si}$/$\operatorname{Ci}$)
integrands where possible and explicitly nested integral representations
where not. After that, the monochromatic case is specialised.

%======================================================================
\section{Exact reformulation (Stage 0)}\label{sec:exact}
%======================================================================

Units $c=1$. The background is $a(t)=a_1 t$ in conformal time $t$ and
$\bar T(t)=\Tref/(a_1t)$, with $\Tref$ an arbitrary constant. The Friedmann
equation $\mathcal H^2=\tfrac{8\pi G}{3}a^2\rho$ fixes
$8\pi G\alpha\Tref^4=3a_1^2$.

\subsection{Velocity potential and temperature frame}

\begin{lemma}[Irrotational radiation fluid]\label{lem:kelvin}
For $P=\rho/3$, $\rho=\alpha T^4$ (zero chemical potential), the relativistic
Euler equation is
$a_\mu\equiv u^\nu\nabla_\nu u_\mu=-h_\mu{}^\nu\partial_\nu\ln T$.
For $w_\mu\equiv Tu_\mu$ this is equivalent to $u^\mu(\dd w)_{\mu\nu}=0$.
Then $\mathcal L_u\dd w=\dd(i_u\dd w)=0$, so the vorticity 2-form $\dd w$ is
Lie-dragged. If $\dd w=0$ on the initial data, then $\dd w=0$ everywhere at
every order, and $Tu_\mu=-\Tref\,\partial_\mu t$ for a scalar $t$, the
velocity-potential time.
\end{lemma}
\noindent
(This is a one-line identity: $u^\mu(\partial_\mu w_\nu-\partial_\nu w_\mu)
=Ta_\nu+h_\nu{}^\mu\partial_\mu T$.) Irrotationality of the initial data is
part of our definition of ``no independent higher-order modes''
(Definition~\ref{def:ic}). It is \emph{not} automatic: the constant-vorticity
mode of radiation is regular as $t\to0$ and must be excluded explicitly.

Take coordinates $(t,X^i)$, where $X^i$ are fluid labels ($u^\mu\partial_\mu
X^i=0$). Because $u_\mu\propto\partial_\mu t$, the shift vanishes and
\begin{equation}\label{eq:metric}
g=N^2\big(-\dd t^2+\gamma_{ij}(t,X)\,\dd X^i\dd X^j\big),\qquad
N=\frac{\Tref}{T_{\rm phys}},\qquad u=N^{-1}\partial_t .
\end{equation}
This is \emph{exact}. We call $\tilde g\equiv g/N^2$ the \emph{temperature
frame}. In it, the fluid worldlines are synchronous and comoving. The
background is $N=a_1t$, $\gamma_{ij}=\delta_{ij}$. Every coordinate here is a
physical label: $t$ is the fluid's velocity potential (normalised by
$t\to0$ as $T\to\infty$), and $X^i$ labels fluid elements. The only residual
freedom is a time-independent relabelling $X\mapsto X'(X)$, which we fix
through the initial data (\S\ref{sec:ic}).

\begin{lemma}[Continuity $\Leftrightarrow$ expansion-free temperature frame]\label{lem:L1}
For the metric \eqref{eq:metric} with arbitrary $N,\gamma_{ij}$ and the fluid
at rest with $\rho\propto N^{-4}$, the spatial components of
$\nabla_\mu T^{\mu\nu}$ vanish identically, i.e.\ Euler holds automatically.
Moreover
$\nabla_\mu T^{\mu}{}_t=-\tfrac23\rho\,\partial_t\ln\det\gamma$.
Hence $\det\gamma_{ij}(t,X)=\det\gamma_{ij}(0,X)$, and
$\sigma_{ij}\equiv\tfrac12\partial_t\gamma_{ij}$ is exactly traceless.
\end{lemma}
\status{verified: \texttt{check\_exact\_lemmas.py}, L1, arbitrary $N$ and six arbitrary $\gamma_{ij}$}

\begin{lemma}[Exact redshift]\label{lem:redshift}
(a) Under $g=N^2\tilde g$, null geodesics are preserved, with
$\tilde p^\mu=N^2p^\mu$, $\tilde u^\mu=Nu^\mu$, and so $E=-g(u,p)=\tilde
E/N$. (b) In the synchronous metric $\tilde g$,
$\tilde\Gamma^t{}_{t\mu}=0$ and $\tilde\Gamma^t{}_{ij}=\sigma_{ij}$, so
$\dd\ln\tilde E/\dd t=-\sigma_{ij}e^ie^j$ with $e^i=\tilde p^i/\tilde p^t$
and $\gamma_{ij}e^ie^j=1$. The ray direction obeys
\begin{equation}\label{eq:ray}
\frac{\dd e^i}{\dd t}=-2\sigma^i{}_je^j-{}^{(\gamma)}\Gamma^i{}_{jk}e^je^k+e^i\sigma_{jk}e^je^k,
\qquad \frac{\dd X^i}{\dd t}=e^i.
\end{equation}
(c) Consequently, with $N_e=\Tref/T_\star$ at emission:
\begin{equation}\label{eq:Tobs}
T_{\rm obs}(\nhat)=T_\star\frac{E_o}{E_e}=T_{\rm fluid}(o)\,
\exp\Big[-\int_{t_e(\nhat)}^{t_o}\sigma_{ij}\big(t,X(t)\big)e^ie^j\,\dd t\Big].
\end{equation}
\end{lemma}
\status{verified: L2, L3}. Isotropic Planck emission in the fluid frame
together with Liouville's theorem makes $T_{\rm obs}$ frequency-independent,
so \eqref{eq:Tobs} is the complete observable.

\begin{remark}[What \eqref{eq:Tobs} buys]
$T_{\rm fluid}(o)$ is direction-independent, so it cancels \emph{exactly} in
$\Theta=T_{\rm obs}/\sky{T_{\rm obs}}-1$, at every order. All
emission-surface, observer and normalisation effects reduce to:
(i) the endpoint $t_e(\nhat)$ fixed by $N(t_e,X(t_e))=\Tref/T_\star$;
(ii) the observation event $t_o$ fixed by proper time;
(iii) the initial direction $e^i(t_o)$ fixed by the tetrad;
(iv) the ray path from \eqref{eq:ray}; and
(v) the sky average.
There is no separate SW, ISW or Doppler term to keep consistent. At first
order we prove the equivalence with the Newtonian-gauge SW+Doppler+ISW
expression for all $\ell$ (\S\ref{sec:first}).
\end{remark}

\subsection{Field equations in temperature-frame variables}
With $n\equiv\ln N$, $D$ and $\Delta$ the $\gamma$-covariant derivative and
Laplacian, a prime for $\partial_t$, and $\sigma^2=\sigma_{ij}\sigma^{ij}$,
Einstein's equations $R_{\mu\nu}=8\pi G\,T_{\mu\nu}$ (the source is
traceless) read \emph{exactly}:
\begin{align}
&-\sigma^2-3n''+\Delta n+2|Dn|^2=\frac{3a_1^2}{N^2},\label{eq:E00}\\
&D_j\sigma^j{}_i-2\partial_in'+2\sigma^j{}_i\partial_jn+2n'\partial_in=0,\label{eq:E0i}\\
&R_{ij}[\gamma]+\sigma_{ij}'-2\sigma_{ik}\sigma^k{}_j-2(D_iD_jn-\sigma_{ij}n')
+\gamma_{ij}(n''-\Delta n)+2D_inD_jn+2\gamma_{ij}(n'^2-|Dn|^2)=\frac{a_1^2}{N^2}\gamma_{ij},\label{eq:Eij}
\end{align}
with the Hamiltonian constraint
$N^2(R[\gamma]-4\Delta n-2|Dn|^2-\sigma^2)+6N'^2=6a_1^2$.
\status{verified as identities: \texttt{check\_exact\_field\_equations.py}
(arbitrary $N$ and a det-1 family with off-diagonal shear) and L4.}

\paragraph{Parametrisation that solves continuity exactly.} Write
\[
N=a_1t\,e^{\nu(t,X)},\qquad
\gamma_{ij}=e^{2\zeta(X)}\big[\exp H(t,X)\big]_{ij},\qquad H_{ii}=0,\qquad H(0,X)=0.
\]
Lemma~\ref{lem:L1} is then satisfied identically, and the unknowns are
$(\nu,H_{ij})$: six functions. At each order $H^{(r)}$ is split into SVT
parts, $H^{(r)}_{ij}=2(\partial_i\partial_j-\tfrac13\delta_{ij}\nabla^2)E^{(r)}
+2\partial_{(i}F^{(r)}_{j)}+h^{(r)}_{ij}$, with $\partial_iF_i=0$ and $h$
transverse-traceless. The induced vector and tensor pieces enter through
$F^{(2,3)}$ and $h^{(2,3)}$.

%======================================================================
\section{Primordial state and initial data}\label{sec:ic}
%======================================================================

\begin{definition}[Gaussian variable and absence of independent higher-order modes]\label{def:ic}
The primordial state is the unique solution of
\eqref{eq:E00}--\eqref{eq:Eij} such that, as $t\to0$:
\begin{enumerate}
\item[(G)] $\gamma_{ij}(t,X)\to e^{2\zeta(X)}\delta_{ij}$, i.e.\ $H\to0$, and
$\nu\to0$, with $\zeta=\eps\,\zeta_1$ \emph{exactly}, where $\zeta_1$ is a
statistically homogeneous and isotropic Gaussian field with power spectrum
$P_\zeta(k)$;
\item[(T)] no primordial tensor modes: the limiting 3-metric is conformally
flat (its Cotton--York tensor vanishes). The labels $X$ are chosen
conformally Cartesian for it, which is what makes (G) possible;
\item[(V)] no vorticity: $\dd(Tu)=0$ (Lemma~\ref{lem:kelvin});
\item[(R)] regularity: $\nu$ and $H$ are analytic in $t^2$ with
$X$-dependent coefficients. This excludes every decaying mode.
\end{enumerate}
``No independent higher-order primordial modes'' means
$\zeta^{(r\ge2)}\equiv0$ in $\zeta=\sum_r\eps^r\zeta^{(r)}/r!$, together with
(T), (V) and (R) at every order.
\end{definition}

\paragraph{Identification of $\zeta$.} At leading order in gradients,
$\nu=O(k^2t^2)$, so uniform-density slices ($N=$const) and $t=$const slices
differ by $\delta t=O(k^2t^3)$. The spatial metric on uniform-density slices
therefore tends to $a^2e^{2\zeta}\delta_{ij}$ with the \emph{same} $\zeta$.
Our $\zeta$ is thus the nonlinear $\delta N$ / uniform-density curvature
perturbation (Lyth--Malik--Sasaki; Langlois--Vernizzi), which is conserved on
super-horizon scales for adiabatic perturbations. The Gaussian variable is
$\zeta$. Consequently $\Phi$, $\delta$ and $v$ in any gauge are non-Gaussian
at second order with specific, derived coefficients.

\paragraph{Residual label freedom.} (G)+(T) fix $X$ up to the conformal group
of $\mathbb R^3$ (translations, rotations, dilations, special conformal
transformations), acting together with the corresponding transformation of
$\zeta$. Translations and rotations are statistical symmetries. Dilations and
SCTs are Weinberg's adiabatic modes, and they supply nontrivial checks (V9,
V10).

\begin{proposition}[Existence and uniqueness order by order]\label{prop:unique}
At order $r$, \eqref{eq:E00}--\eqref{eq:Eij} take the form
$\mathcal L[\nu^{(r)},H^{(r)}]=S^{(r)}[\text{lower orders}]$, where
$\mathcal L$ is the first-order operator. The regular homogeneous solutions
of $\mathcal L$ are: the scalar growing mode (equivalent to adding
$\zeta^{(r)}$), the regular tensor mode $h\propto j_0(kt)$, and
$\partial_tF=0$ (a relabelling). (G), (T) and the label choice set all three
to zero. The particular solution is given by retarded Green's functions with
lower limit $t=0$; the integrals converge because $S^{(r)}$ is regular by
induction. So the solution exists and is unique.
\end{proposition}
\status{Stage 2: the regular-homogeneous-solution classification is to be
proved from the explicit linear operator. The scalar part is already
confirmed at first order.}

\paragraph{Linear operator and Green's functions.} With the
comoving-curvature master variable $\mathcal R=\nu+\zeta-\tfrac13\nabla^2E$:
\[
\Big(\partial_t^2+\tfrac2t\partial_t-\tfrac13\nabla^2\Big)\mathcal R=S_{\mathcal R},\qquad
\Big(\partial_t^2+\tfrac2t\partial_t-\nabla^2\Big)h_{ij}=S^{TT}_{ij},\qquad
\partial_tF_i=S^{V}_i .
\]
The regular retarded Green's functions are
$G_{\mathcal R}(t,t')=\frac{t'}{t}\frac{\sqrt3}{k}\sin\frac{k(t-t')}{\sqrt3}$
and $G_h(t,t')=\frac{t'}{t}\frac{\sin k(t-t')}{k}$. The remaining scalar
fields follow algebraically or by one time integral from the constraints.
\status{the operator is confirmed at first order; the extraction of the
$S^{(2,3)}$ sources is Stage 2, verified by substituting back into the
complete Einstein tensor (V3).}

\paragraph{Initial data at finite $t_i$.} The problem asks for initial data
that satisfies the constraints. It is obtained by evaluating the regular
solution at $t_i$; equivalently, it is the gradient-expansion series in
$t_i^2$. The Hamiltonian and momentum constraints hold identically at every
order (V4). The higher-order metric and fluid perturbations are \emph{not}
zero at $t_i$. For example, $\gamma^{(2)}_{ij}(t\to0)=4\zeta_1^2\delta_{ij}$,
which is the $\eps^2/2$ coefficient of $e^{2\eps\zeta_1}$, and $\nu^{(2)}$
and $H^{(2)}$ start at $O(t^2)$ with coefficients fixed by \eqref{eq:E0i} and
the Hamiltonian constraint. The limit $t_i\to0$ is shown to commute with the
observable because all kernels are continuous at $t_i=0$.

%======================================================================
\section{The observable, exactly, then order by order}\label{sec:obs}
%======================================================================

\paragraph{Physical measurement labels.} Every one of these is a scalar or a
physically defined frame:
\begin{itemize}
\item Observer: the fluid line $X=X_o$.
\item Initial event $\mathcal I$: the point on that line where
$T_{\rm phys}=T_I$, so $N(t_I,X_o)=\Tref/T_I$. The limit $T_I\to\infty$,
which measures from the singularity, exists by (R).
\item Observation event: $\tau_o=\int_{t_I}^{t_o}N(t,X_o)\,\dd t$ at fixed
$\tau_o$. This defines a \emph{stochastic} $t_o=\bar t_o+\eps\,\delta t_o+\dots$,
with $\bar t_o^2=\bar t_I^2+2\tau_o/a_1$.
\item Emission event: the first intersection of the past-directed ray with
$N=\Tref/T_\star$, giving $t_e(\nhat)=\bar t_\star+\dots$ with
$\bar t_\star=\Tref/(a_1T_\star)$.
\item Tetrad: Fermi--Walker transported along $u$ in $g$. This is equivalent
to parallel transport along the geodesic $\tilde u$ in $\tilde g$ (FW
transport of vectors orthogonal to $u$ is conformally covariant, a Stage 0
item to be written up). Then $\dd E_a{}^i/\dd t=-\sigma^i{}_jE_a{}^j$, with
initial orientation $E_a{}^i(t_I)=[\gamma(t_I,X_o)^{-1/2}]_{ia}$ (symmetric
square root, aligned with the conformally Cartesian labels). The ray then
starts from $e^i(t_o)=-E_a{}^i(t_o)\nhat^a$.
\end{itemize}

\begin{proposition}[Orientation independence of $C_\ell$]
Two admissible tetrads at the observation event differ by an $SO(3)$
rotation, possibly realisation-dependent. $\sum_m|a_{\ell m}|^2$ is invariant
under rotations of the map \emph{realisation by realisation}. So $C_\ell$ is
independent of the initial orientation and of the transport law. Only the
non-rotational (symmetric) part of $E_a{}^i$, which orthonormality fixes,
affects $C_\ell$.
\end{proposition}
\noindent
We nevertheless carry the specified tetrad. The independence then serves as
check V12.

\paragraph{Normalisation (exact through third order).} Write
$T_{\rm obs}=T_{\rm fluid}(o)e^{-I(\nhat)}$ with
$I=\eps I_1+\tfrac{\eps^2}2I_2+\tfrac{\eps^3}6I_3$, and let $\sky{\cdot}$ be
the sky average. This is a realisation-dependent \emph{random} number, not an
ensemble average. Then (\texttt{normalization\_expansion.py}):
\begin{align*}
\Theta^{(1)}&=\sky{I_1}-I_1,\\
\Theta^{(2)}&=I_1^2-I_2-2\sky{I_1}I_1+2\sky{I_1}^2-\sky{I_1^2}+\sky{I_2},\\
\Theta^{(3)}&=-I_3+3I_1I_2-I_1^3+3\sky{I_1}I_1^2-3\sky{I_1}I_2-3\sky{I_2}I_1+3\sky{I_1^2}I_1
-6\sky{I_1}^2I_1\\&\quad+\sky{I_3}-3\sky{I_1I_2}+\sky{I_1^3}+6\sky{I_1}^3-6\sky{I_1}\sky{I_1^2}+6\sky{I_1}\sky{I_2},
\end{align*}
with $\sky{\Theta^{(r)}}=0$ identically (verified). Products such as
$\sky{I_1}I_1$ survive in $a_{\ell\ge1,m}$ and contribute to both
$C^{(22)}$ and $C^{(13)}$. They are the ``normalisation interference''
terms.

\paragraph{Content of $I_1,I_2,I_3$.} Expand \eqref{eq:Tobs} about the
background ray $\bar X(t)=X_o+\nhat(\bar t_o-t)$. Every term is kept:
\begin{itemize}
\item $I_1=\int_{\bar t_\star}^{\bar t_o}\sigma^{(1)}_{ij}\nhat^i\nhat^j\,\dd t$.
\item $I_2$: $\sigma^{(2)}$ on the background ray; the deflection term
$2\sigma^{(1)}_{ij}\nhat^i\delta e^j$ with $\delta e$ from \eqref{eq:ray}; the
Born correction $\delta X\cdot\partial\sigma^{(1)}$; the endpoint terms
$\delta t_e\,\sigma^{(1)}|_e$ and $\delta t_o\,\sigma^{(1)}|_o$ (plus the
$\delta t_o$ shift of the whole ray); and the tetrad term from
$E^{(1)}=-\tfrac12\gamma^{(1)}$ at the observer, which acts as aberration by
the shear.
\item $I_3$: all of the above one order higher, including post-Born
(second-order deflection, mixed $\delta X^{(1)}\delta e^{(1)}$, and
$\tfrac12\delta X\delta X\partial\partial\sigma^{(1)}$), $\delta t_e^{(2)}$
and $\delta t_o^{(2)}$ from the nonlinear endpoint conditions, products of
endpoint and deflection shifts, and the second-order tetrad.
\end{itemize}
The endpoint conditions are solved perturbatively and exactly. For example,
$\nu(t_e,X(t_e))+\ln(t_e/\bar t_\star)=0$ gives
$\delta t_e^{(1)}=-\bar t_\star\nu^{(1)}(\bar t_\star,\bar X_\star)$, and
higher orders follow by Lagrange inversion. The null condition
$\gamma_{ij}e^ie^j=1$ is preserved by \eqref{eq:ray}; it is checked at each
order (V6).

\paragraph{Dipole.}\label{sec:dipole} The measured dipole is subtracted
after $\Theta$ has been constructed and normalised, as required. This only
removes $a_{1m}$. All couplings of the kinematic, observer-induced and
intrinsic dipole to other multipoles are already inside $\Theta^{(2,3)}$:
they come from the tetrad, the observer endpoint and the normalisation
products above. So $C_{\ell\ge2}$ is unaffected by the subtraction, while
$C_0\equiv0$ and $C_1\equiv0$ by construction.

%======================================================================
\section{First order: complete and verified (Stage 1)}\label{sec:first}
%======================================================================

With $x\equiv kt/\sqrt3$ and one mode $\zeta_k e^{ik\cdot X}$, the regular
solution is
\[
\nu=-\zeta\Big(\frac{\sin x}x-\cos x\Big),\qquad
\gamma^{(1)}_{ij}=2\Big[\zeta-\tfrac{1}{3}\nabla^2E\Big]\delta_{ij}+2\partial_i\partial_jE,\qquad
k^2E=3\zeta\,s(x)-3\zeta,\quad s(x)\equiv\frac{2\sin x}x-\cos x,
\]
and $\mathcal R=\zeta\sin x/x$.
\status{verified: all ten Einstein equations, regularity, the $t\to0$ limits
and expansion-freeness (\texttt{check\_first\_order.py}).} The
Newtonian-gauge solution $\Phi=\Psi=-2\zeta(\sin x-x\cos x)/x^3$, with
$\delta_N$ and $v_N$, is verified independently in the same script.

With the Fourier convention
$\zeta(X)=\int\!\frac{\dd^3k}{(2\pi)^3}e^{ik\cdot X}\zeta_k$ and
$\ens{\zeta_k\zeta_{k'}}=(2\pi)^3\delta^3(k+k')P_\zeta(k)$,
\[
\Theta^{(1)}(\nhat)=\int\!\frac{\dd^3k}{(2\pi)^3}\zeta_ke^{ik\cdot X_o}\sum_{\ell\ge1}(2\ell+1)i^\ell\,\mathcal T_\ell(k)P_\ell(\hat k\cdot\nhat),\qquad
\mathcal T_\ell=-3\!\int_{x_\star}^{x_o}\!\dd x\;s'(x)\,w_\ell\big(\sqrt3(x_o-x)\big),
\]
where $w_\ell(y)=\tfrac13j_\ell(y)+j_\ell''(y)$, $x_\star=k\bar t_\star/\sqrt3$
and $x_o=k\bar t_o/\sqrt3$. Invariantly,
$x=k/(\sqrt3\,aH)$ at the corresponding event, and $x_o/x_\star=T_\star/\bar
T_o$. The factor $w_\ell$ is the Legendre image of the traceless projector
$\tfrac13-(\hat k\cdot\nhat)^2$; the redshift is pure shear.

\paragraph{Gauge-invariance proof at first order (all $\ell$).} The
Newtonian-gauge expression is
$\mathcal T^N_\ell=[\tfrac14\delta_N+\Phi]_\star j_\ell-kv_N(t_\star)j_\ell'
+\int2\Phi'j_\ell+kv_N(t_o)j_\ell'(0)$: emission-surface temperature shift,
SW, emitter and observer Doppler, and ISW. We show
$\partial_{t_\star}(\mathcal T_\ell-\mathcal T^N_\ell)=0$ \emph{for symbolic
$\ell$}, using only the Bessel ODE, and that both vanish at $t_\star=t_o$ for
$\ell\ge1$. Hence $\mathcal T_\ell\equiv\mathcal T^N_\ell$.
\status{verified}

\paragraph{Closed form.} Rayleigh's formula makes the integrand a finite sum
of $e^{i\beta u}\times$(rational in $u=x_o-x$), with
$\beta\in\{\pm1,\pm\sqrt3,\pm1\pm\sqrt3\}$. Term-by-term integration gives
$\mathcal T_\ell$ as an elementary combination of $\sin$, $\cos$, powers,
and $\operatorname{Si}$, $\operatorname{Ci}$ at arguments
$(\sqrt3\pm1)(x_o-x_\star)$, $(\sqrt3\pm1)x_o$, $(\sqrt3\pm1)x_\star$, with
cancelling logarithms at the observer endpoint
(\texttt{transfer\_closed\_form.py}). Each antiderivative is verified by exact
differentiation, and the cancellation of the observer-endpoint
singularities is asserted.

\paragraph{Spectrum.} The angular normalisation is fixed by
$\ens{a^{(1)}_{\ell m}a^{(1)*}_{\ell m}}=(4\pi)^2\!\int\!\frac{\dd^3k}{(2\pi)^3}P_\zeta|\mathcal T_\ell|^2|Y_{\ell m}(\hat k)|^2$, which gives
\[
C^{(11)}_\ell=4\pi\int\frac{\dd k}{k}\Delta^2_\zeta(k)\,\mathcal T_\ell(k)^2,\qquad
\Delta^2_\zeta=\frac{k^3P_\zeta}{2\pi^2}.
\]
For $\Delta^2_\zeta(k)=A_\zeta\,\delta_D(\ln(k/k_0))$ this is
$C^{(11)}_\ell=4\pi A_\zeta\,\mathcal T_\ell(k_0)^2$. Here
$A_\zeta=\ens{\zeta_1^2}$ is the quantity held fixed as $k_0$ varies.

%======================================================================
\section{Second and third order (Stages 2--3)}\label{sec:higher}
%======================================================================

\paragraph{Kernels.} For each order-$r$ field $Y$,
\[
Y^{(r)}(t,X)=\int\prod_{a=1}^r\frac{\dd^3k_a}{(2\pi)^3}\zeta_{k_a}\;e^{i(\sum k_a)\cdot X}\;\mathcal K^Y_r(k_1,\dots,k_r;t),
\]
with symmetrised kernels. At $r=2$ the sources are products of
trigonometric functions of $k_at/\sqrt3$ and rational functions of $t$, and
$G_{\mathcal R}$ and $G_h$ are elementary. So $\mathcal K_2$ is elementary
plus $\operatorname{Si}/\operatorname{Ci}$ of $(\pm k_1\pm k_2\pm\sqrt3k)t/\sqrt3$
and $(\pm k_1\pm k_2\pm k)t/\sqrt3$, the same structure as the induced-GW
kernel in radiation domination (Kohri--Terada) but at finite time and for
all sectors. The vector part $F^{(2)}$ is nonzero even though the vorticity
is zero; it is a metric frame-dragging term fixed by \eqref{eq:E0i}. At
$r=3$ the sources contain $\mathcal K_2$, so $\mathcal K_3$ is elementary
plus one-dimensional definite integrals of
$\operatorname{Si}/\operatorname{Ci}\times$(elementary). These are given as
explicit nested integrals over $t'\in[0,t]$ with complete integrands. We do
not expect general closed forms, and parts that do close (e.g.\
$\int\operatorname{Si}(\alpha s)\sin(\beta s)s^{-m}\dd s$ in terms of
$\operatorname{Si}/\operatorname{Ci}$ and ${}_2F_3$) will be evaluated where
they appear.

\paragraph{Sky projection.} Every term in $I_r$ is a product of
$e^{ik_a\cdot\nhat\,r_a}$ (different line-of-sight distances $r_a$ for
nested integrals), polynomials in $\hat k_a\cdot\nhat$, and sky averages.
Polynomials are generated by $\partial_{r_a}$. The identity
$\int\dd\Omega\,Y^*_{\ell m}(\nhat)e^{iK\cdot\nhat}=4\pi i^\ell
j_\ell(|K|)Y^*_{\ell m}(\hat K)$ with $K=\sum_ak_ar_a$, together with
$\sum_mY_{\ell m}(\hat K)Y^*_{\ell m}(\hat K')=\frac{2\ell+1}{4\pi}P_\ell(\hat
K\cdot\hat K')$, makes every $C^{(rs)}_\ell$ an explicit integral. No
flat-sky, Limber or Born approximation is used.

\paragraph{Wick contractions (Gaussian $\zeta_1$).}
\begin{align*}
C^{(22)}_\ell&=\frac{2}{2\ell+1}\sum_m\int\!\frac{\dd^3k_1\dd^3k_2}{(2\pi)^6}P(k_1)P(k_2)\,\big|\mathcal A^{(2)}_{\ell m}(k_1,k_2)\big|^2\quad(\text{connected}),\\
C^{(13)}_\ell&=\frac{3}{2\ell+1}\sum_m\int\!\frac{\dd^3k}{(2\pi)^3}P(k)\,\mathcal A^{(1)}_{\ell m}(k)\int\!\frac{\dd^3q}{(2\pi)^3}P(q)\,\mathcal A^{(3)*}_{\ell m}(k,q,-q),
\end{align*}
where $\mathcal A^{(r)}_{\ell m}$ is the $a_{\ell m}$ kernel; the factors 2
and 3 count the pairings of symmetrised kernels. For $\ell\ge1$ the
disconnected part $\ens{a^{(2)}}\ens{a^{(2)*}}$ vanishes by isotropy, and
$a_{00}\equiv0$ by normalisation. The sky-average terms
($\sky{I_1}I_1$, etc.) are included in $\mathcal A^{(2,3)}$ as
$\ell'=0$ projections inside the same Fourier integrals.

\paragraph{Monochromatic specialisation (Stage 5).} Insert
$P(k)=\frac{2\pi^2A_\zeta}{k_0^3}\delta_D(k-k_0)$, i.e.\
$\Delta^2_\zeta=A_\zeta\delta_D(\ln k/k_0)$, \emph{after} the complete
observable is built. In $C^{(22)}$, rotational invariance leaves the single
angle $\mu=\hat k_1\cdot\hat k_2\in[-1,1]$ plus the internal line-of-sight
times; $|k_1+k_2|=k_0\sqrt{2+2\mu}$ covers $[0,2k_0]$. In $C^{(13)}$ only
$\mu=\hat k\cdot\hat q$ survives. Held fixed: $A_\zeta=\ens{\zeta_1^2}$.
Results depend on $k_0$ only through $x_\star(k_0)$, $x_o(k_0)$ and $x_I(k_0)$.

\paragraph{Exchange of limits.} At finite $t_\star,t_o$ every kernel is a
continuous, bounded function of the momenta on the compact sets selected by
the $\delta$-function spectrum. Pairing the distribution with it is
therefore legitimate, and Fubini applies to the finite nested time
integrals. The $t_I\to0$ and $t_i\to0$ limits commute with all integrals by
dominated convergence, since the kernels are analytic in $t^2$. The
\emph{late-time} limit $x_o\to\infty$ is \emph{not} uniform near the
sound-speed resonance $|k_1|+|k_2|=\sqrt3|k_1+k_2|$, i.e.\
$\mu=-1/3$ for the monochromatic case (tensor sector). There, logarithms like those of the
radiation-era induced-GW kernel appear. This limit is taken only at the end,
with the $\mu$-integral done first, and is flagged in \S\ref{sec:open}.

%======================================================================
\section{Validation matrix}\label{sec:val}
%======================================================================

All checks are exact symbolic identities. Numerical spot checks at 30
digits appear only as \emph{supplementary} sanity output and are never used
as evidence.

\begin{longtable}{@{}p{0.6cm}p{6.4cm}p{6.0cm}p{1.7cm}@{}}
\toprule
& Check & Method & Status\\\midrule\endhead
V1 & Exact lemmas L1--L4, field equations \eqref{eq:E00}--\eqref{eq:Eij} & identities for arbitrary metric functions & done\\
V2 & Background & substitution & done\\
V3 & Order-$r$ solution satisfies all 10 Einstein equations & substitute kernels into the \emph{full} $G_{\mu\nu}$ of \eqref{eq:metric}, expanded to $O(\eps^r)$; Si/Ci derivatives are elementary; nested integrals via Leibniz and Green's identity & $r{=}1$ done; $r{=}2,3$ planned\\
V4 & Constraints on initial data at finite $t_i$; regularity; $t\to0$ limits (G),(T) & series in $t$ & $r{=}1$ done\\
V5 & Fluid normalisation, continuity and Euler & exact by Lemma~\ref{lem:L1}; re-checked per order & done/each order\\
V6 & Photon null condition and geodesic equation for perturbed rays & residual of \eqref{eq:ray} and of $\tilde g(\tilde p,\tilde p)$ per order & planned\\
V7 & Gauge invariance: temperature frame vs Newtonian/Poisson gauge & full independent recomputation in Poisson gauge, compared multipole by multipole & $r{=}1$ done (all $\ell$); $r{=}2$ planned\\
V8 & Gauge transformation of the solution & transform $\mathcal K_2$ to Poisson gauge; check Einstein equations there and match the radiation-era second-order initial conditions in the literature & planned\\
V9 & Constant-$\zeta$ (dilation) Ward identity & $\Theta[\zeta_s+c](\nhat)=\Theta[\zeta_s(e^{-c}\cdot)](\nhat)$ $\Rightarrow$ squeezed limits of $\mathcal K_2,\mathcal K_3$ & $r{=}1$ done; planned\\
V10 & Gradient-$\zeta$ (SCT) Ward identity & $O(q)$ squeezed terms cancel in $\Theta$ & planned\\
V11 & Normalisation: $\sky{\Theta^{(r)}}=0$ & algebraic & done\\
V12 & Tetrad-orientation independence of $C_\ell$ & add an arbitrary $\eps$-dependent rotation & planned\\
V13 & Coincidence limit $t_\star\to t_o$: $\Theta^{(r)}\to0$ & series in $t_o-t_\star$ & $r{=}1$ done\\
V14 & Super-horizon limit $k\to0$: no $\ell\ge2$ anisotropy at $O(k^0)$ & series & $r{=}1$ done\\
V15 & Antiderivatives of closed forms & exact differentiation & done ($\mathcal T_\ell$)\\
V16 & Wick and angular normalisation & two routes: Fourier and configuration-space correlation $\ens{\Theta(\nhat)\Theta(\nhat')}$ & planned\\
V17 & Independent CAS & re-derive key identities in Maxima (installed by the Containerfile) & planned\\
V18 & Independent agents & blind re-derivation and audit & see \S\ref{sec:indep}\\
\bottomrule
\end{longtable}

%======================================================================
\section{Plan, stages and cost}\label{sec:plan}
%======================================================================
\begin{enumerate}
\item[S0] Exact framework. \emph{Done} apart from the write-up of the
FW-conformal lemma.
\item[S1] First order. \emph{Done}: solution, gauge equivalence, closed-form
$\mathcal T_\ell$, $C^{(11)}_\ell$.
\item[S2] Second-order kernels $(\nu,E,F,h)^{(2)}$: extract the sources from
\eqref{eq:E00}--\eqref{eq:Eij} by exact expansion, solve with
$G_{\mathcal R}$ and $G_h$, and run V3, V4, V8, V9. This also yields the
Poisson-gauge $\Phi^{(2)},\Psi^{(2)},\omega^{(2)},h^{(2)}$ as a
by-product.
\item[S3] Third-order kernels as nested integrals, with V3 by
differentiation under the integral sign. This is the costliest symbolic
step; use SymEngine or FORM for the tensor expansion.
\item[S4] Observable: $I_2$ and $I_3$ with all endpoint, observer, tetrad
and normalisation terms (V6, V11--V14), $C^{(22)}$ and $C^{(13)}$ (V16), and
the dilation/SCT checks on the full $C_\ell$ (V9, V10).
\item[S5] Monochromatic specialisation and asymptotics ($x_o\to\infty$,
$x_\star\ll1$, $x_\star\gg1$), with a justification for each exchange of
limits.
\end{enumerate}

%======================================================================
\section{Open items and risks}\label{sec:open}
%======================================================================
\begin{enumerate}
\item No closed form is expected for the generic third-order kernel. It
will be given as explicit nested integrals, and parts that close will be
marked as \emph{evaluated}.
\item The late-time limit is non-uniform at the sound-speed resonance
$\mu=-1/3$. The order of limits must be stated, and only the $\mu$-first
order is justified.
\item IR behaviour for non-monochromatic spectra: $C^{(13)}$ and $C^{(22)}$
are separately sensitive to $q\to0$. The dilation and SCT identities
(V9, V10) predict cancellation up to a physical rescaling of the short-mode
spectrum. This must be shown, not assumed. It does not arise for the
monochromatic spectrum.
\item The perturbative uniqueness of the emission intersection and the
absence of caustics hold order by order, but this is a statement about the
asymptotic series and not about convergence in $\eps$.
\item Proposition~\ref{prop:unique} still requires the explicit
classification of regular homogeneous solutions in the vector and tensor
sectors.
\end{enumerate}

\section{Independent validation by subagents}\label{sec:indep}
Two separately prompted agents reviewed this proposal and the scripts. One
was a physics referee; the other did a blind first-order re-derivation and
an audit of the scripts. Their findings, and how they were resolved, are
recorded in \texttt{approach/VALIDATION.md}.

\end{document}
